\section{Benchmark adapters and coverage}
\label{sec:benchdetails}

\paragraph{LongMemEval-V2 haystacks.}
On the \texttt{small} tier each domain has exactly one haystack, identical
across questions even as an ordered tuple, so every memory is available to
every question in its domain and there is no availability confound between
arms. The same test returns 222 and 211 distinct haystacks on
\texttt{medium}, a positive control showing that it distinguishes the two
cases. Trajectory lengths were measured with the benchmark's own tokenizer
over 200 trajectories and 5{,}095 states.

\paragraph{Scorer dispatch.}
There are eight distinct evaluation-specification strings, not six: two
phrase-match variants differ only in their separator. The adapter therefore
dispatches on the full string, never on the function name, and calls the
benchmark's scorer verbatim. Its phrase matcher tests gold-phrase
\emph{containment}, not set equality, which a name-based reimplementation
gets wrong. An LLM-judged task raises an error rather than returning a
default; a silent \texttt{False} would turn ``cannot evaluate'' into ``wrong
answer'' for 156 of 451 questions in both arms.

\paragraph{Pairing and power.}
Seeded deterministic pairing within domain, stratified by the benchmark's own
question-type taxonomy, gives 1{,}889 memories $\times$ 8 tasks $=$ 15{,}112
cells, and every memory was verified to receive exactly 8. The build asserts
that the same seed reproduces the grid byte-identically \emph{and} that a
different seed changes it, so the reproducibility check is not vacuous. That
grid is the available pairing, not the executed design. The floor of 60
memories was fixed before running from a recovery study on synthetic data
with known planted components; the same study showed that widening the task
axis helps more than adding memories.

\paragraph{BEAM.}
BEAM~\cite{beam} was in our original plan as a multi-domain check but cannot
supply the task axis: 92.6\% of its cited turns are cited by exactly one
question. Being CC~BY-SA~4.0, it would also have to ship separately, since
ShareAlike would otherwise relicense the rest of the release.

\paragraph{SciCode.}
The memory unit is a reusable helper routine with its procedure, which
transfers across problems. Against an early feasibility estimate of 8--12
pairs, the full candidate pool measured on the built adapter is $\sciPairs$
donor$\to$recipient pairs ($\sciMems$ memories, $\sciTasks$ tasks), and we
ran all of it.

\paragraph{SWE-bench Verified.}
Pairing is plentiful but the \emph{measurable} yield is low. Across
$\covSwRepos$ repositories, after a localizer that misses the edit region on
$\sweLocMiss$ cells and edits that neither apply nor produce a test verdict,
$\sweMatched$ cells are interpretable under all three representations, and
$\sweDomRepoPct\%$ come from one repository. A scoring defect was corrected before these counts: a test suite
that printed no verdict at all (an import or collection error) had been
recorded as a reward of zero. Those $\lfSweDemoted$ arm-reads
($\lfSweDemDjango$ django, $\lfSweDemSympy$ sympy) are now unscorable, and
every contrast includes zero after the correction. The matched set of
$\lfSweMatched$ cells is pooled from four runs across the two repositories; it grew as runs were added, not through the correction, which can only
remove cells. Table~\ref*{M-tab:coverage} of the main paper counts SWE-bench in
(memory, task, representation) cells; counts per recipient task (107/37/70)
describe a different unit.
